% 08_related_work.tex — synthesises literature/related_work.md.

\section{Related Work}\label{sec:related}

We situate our work in six clusters of prior literature,
organised by the boundary or methodology they target.

\subsection{LLM tool-use and prompt injection}

Greshake et al.~(2023) introduced the concept of \emph{indirect
prompt injection} via tool and retrieval content, demonstrating
that LLM-integrated applications are broadly exploitable
through content they fetch~\cite{greshake_2023_indirect}. Zou
et al.~(2023) developed automated, transferable adversarial
suffixes for aligned LLMs via greedy coordinate-gradient
search~\cite{zou_2023_universal}. Chowdhury et al.~(2024)
provided a comparative survey of attacks on LLMs and showed
that defenses fail in isolation and under
composition~\cite{chowdhury_2024_defenses} --- the same
empirical pattern we observe at the MCP layer. Wang et
al.~(2024) built ToolCommander, a framework for LLM
tool-scheduling hijack via injected tool descriptors, which
matches our \texttt{A-NSP-001} and \texttt{A-TOL-005}
findings~\cite{wang_2024_toolcommander}. These works establish
that prompt injection is a structural property of LLM agents
and that tool scheduling is exploitable; we extend both
findings to the MCP protocol layer with concrete per-boundary
attack patterns.

\subsection{MCP and agent-protocol security}

Chen et al.~(2026) conducted the largest empirical study of
runtime MCP servers to date, finding that scanners under-
report runtime risks and that production servers routinely
share stdio workers and follow symlinks across tenant
roots~\cite{chen_2026_mcp_security}. Their work validates our
Phase-1 threat model at scale. An et al.~(2026) introduced
FlowGuard, a behavioural-trace correlation detector for MCP
security~\cite{an_2026_flowguard}; their methodological
approach informs our Evaluator module (\S\ref{sec:framework}).
Lee et al.~(2026) demonstrated mid-session tool injection in
WebMCP, showing that third-party-script registry mutation
bypasses pre-session policies~\cite{lee_2026_webmcp}; we add
\texttt{notifications/tools/list\_changed} audit-trail support
to our reference secure server in response. Liu et al.~(2026)
introduced MCPEvol-Bench, a dynamic-evolution benchmark that
shows static benchmarks overstate robustness~\cite{liu_2026_mcpevol};
mid-session mutations are now first-class in our attack
library. Jing et al.~(2026) argued for \emph{isolation as a
first-class principle} for LLM-agent
safety~\cite{jing_2026_isolation} --- the exact framing of
our paper --- but provided neither a concrete boundary
catalogue nor empirical measurement.

\subsection{Memory and cache attacks}

Torres et al.~(2026) demonstrated long-term memory poisoning
in LLM agents, showing that memory entries persist across
sessions and that write-time detection is
preferred~\cite{torres_2026_memory}. Their findings bridge the
prompt-injection cluster and our memory/cache boundaries: the
agent-side counterpart of our \texttt{A-CCH-001} /
\texttt{A-MEM-001} MCP-side failure modes.

\subsection{Sandboxing and physical isolation}

Dipta et al.~(2024) demonstrated that side channels survive
logical sandboxing: their dynamic frequency-based
fingerprinting attack defeats gVisor, Firecracker, SGX, and
SEV~\cite{dipta_2024_sandbox}. This is a threat-model
limitation we honestly disclose (\S\ref{sec:discussion}); our
work is at the MCP protocol layer, not the silicon layer.
Amazon's Firecracker microVM~\cite{firecracker_github} is the
substrate for our reference-server deployment.

\subsection{Standards and architectural priors}

NIST SP 800-207 (Zero Trust Architecture~\cite{nist_sp_800_207})
provides the architectural prior for our defense blueprint:
\emph{never trust, always verify}. The OWASP LLM Top 10
v1.1~\cite{owasp_llm_top10_v11} provides the LLM-application
taxonomy; we extend it to the protocol layer with explicit
MCP-boundary bindings. The MCP specification and Anthropic's
announcement~\cite{mcp_spec,anthropic_mcp_announce} are the
specification under study; our defense blueprint and spec-
change recommendations target the next revision (MCP-2).

\subsection{LSP analogue}

Zhu et al.~(2025) introduced LSPFuzz, a grammar-aware fuzzing
campaign over $\sim$300 LSP servers that found many
bugs~\cite{zhu_2025_lspfuzz}. LSP and MCP are both JSON-RPC-
based protocols between a host and a language/tool-specific
server; LSPFuzz is the closest architectural analogue and the
methodological template for our \texttt{GrammarFuzzAttack}
(\texttt{A-FUZZ-001}). We adapt the campaign to MCP's eight-
boundary surface.

\subsection{Positioning summary}

\begin{table}[t]
\caption{Positioning vs.\ prior work. \checkmark{} = covered;
$\triangle$ = partial; $\times$ = not covered.}
\label{tab:positioning}
\centering
\small
\begin{tabularx}{\linewidth}{l c c c c c}
\toprule
                                & \rotatebox{70}{Greshake 2023} & \rotatebox{70}{Wang 2024} & \rotatebox{70}{Torres 2026} & \rotatebox{70}{Chen 2026} & \rotatebox{70}{This work} \\
\midrule
Eight-boundary catalogue      & $\times$  & $\times$ & $\times$ & $\triangle$ & \checkmark \\
Concrete per-boundary attacks  & $\times$  & $\triangle$ & $\triangle$ & $\triangle$ & \checkmark \\
Reproducible harness           & $\times$  & $\checkmark$ & $\times$ & $\triangle$ & \checkmark \\
Multi-tenant leakage metric   & $\times$  & $\times$ & $\times$ & $\times$ & \checkmark \\
Defense-composition study      & $\times$  & $\times$ & $\times$ & $\times$ & \checkmark \\
Open attack library            & $\times$  & $\times$ & $\times$ & $\times$ & \checkmark \\
\bottomrule
\end{tabularx}
\end{table}

The novelty score is 8/10 per our pre-registered Phase-4
assessment (\texttt{docs/01\_Research\_Gap.md}); the
empirical numbers in this paper upgrade the conceptual
scaffolding to executed evidence.

\subsection{How our threat model differs}

Our threat model differs from each prior cluster in a specific
way:

\begin{itemize}
  \item \textbf{vs.\ Greshake et al.~(2023)}: their threat
        model is the LLM-integrated application; ours is the
        MCP server. Indirect prompt injection is the common
        surface; the defense locus differs.
  \item \textbf{vs.\ Wang et al.~(2024)}: their threat model
        is the LLM tool scheduler; ours is the MCP namespace
        boundary. ToolCommander demonstrates tool-scheduling
        hijack; our \texttt{A-NSP-001} row is the MCP-specific
        analogue.
  \item \textbf{vs.\ Torres et al.~(2026)}: their threat model
        is agent-side memory; ours is server-side cache +
        memory. The agent-server symmetry is what makes the
        two papers complementary: Torres covers the agent
        half of the boundary; we cover the server half.
  \item \textbf{vs.\ Chen et al.~(2026)}: their threat model
        is the MCP server's single-tenant default behaviour;
        ours is the multi-tenant isolation surface. Chen
        measures single-tenant runtime risk; we measure
        cross-tenant leakage.
  \item \textbf{vs.\ An et al.~(2026)}: their threat model is
        MCP-server signal traces; ours is MCP-server
        protocol-level leakage. FlowGuard detects; we measure.
  \item \textbf{vs.\ Lee et al.~(2026)}: their threat model
        is WebMCP mid-session mutation; ours is MCP static
        surface plus mid-session defence audit. The two
        threat models are bridged by the
        \texttt{notifications/tools/list\_changed} audit-trail
        recommendation in our defense blueprint.
  \item \textbf{vs.\ Liu et al.~(2026)}: their threat model is
        MCP-server tool evolution; ours is MCP-server
        isolation. MCPEvol-Bench studies robustness under
        change; we study leakage under concurrency.
  \item \textbf{vs.\ Jing et al.~(2026)}: their threat model
        is LLM-agent safety in general; ours is MCP
        multi-tenant safety specifically. Jing argues for
        isolation as a first-class principle; we
        operationalise it with concrete boundaries and
        defenses.
  \item \textbf{vs.\ Dipta et al.~(2024)}: their threat model
        is OS-level sandboxing; ours is protocol-level
        isolation. The two papers are complementary: Dipta
        covers what protocol-layer defenses cannot close;
        we cover what OS-layer sandboxes cannot close.
\end{itemize}

The common thread is that \emph{no prior work} combines all
eight boundaries, all six STRIDE letters, a reproducible
multi-tenant harness, a leakage metric, and a defense-
composition study. The combination is the unit of novelty.

\subsection{Cluster F revisited: why LSPFuzz is the closest analogue}

LSPFuzz~\cite{zhu_2025_lspfuzz} is the closest architectural
analogue to MCP because LSP and MCP share three structural
properties:

\begin{enumerate}
  \item Both are JSON-RPC-2.0-based protocols.
  \item Both mediate between a host (an editor or an agent
        runtime) and a language/tool-specific server.
  \item Both are stateless at the protocol level but stateful
        at the server level.
\end{enumerate}

LSPFuzz's grammar-aware fuzzing campaign over $\sim$300 LSP
servers is therefore the methodological template for our
\texttt{GrammarFuzzAttack}; the differences are in the
attack-surface derivation (LSP's surface is grammar-shaped;
MCP's surface is boundary-shaped). We adapt LSPFuzz's
approach to the eight-boundary surface and add the
\texttt{A-CCH-*} cache attacks that have no LSP analogue.

\subsection{Cluster E revisited: standards as the design prior}

Our defense blueprint is consistent with three standards-
level priors:

\begin{itemize}
  \item \textbf{NIST SP 800-207 (Zero Trust)}
        \cite{nist_sp_800_207}: ``never trust, always verify''
        is the design principle behind all four defenses.
        Each defense re-validates the tenant identity at a
        different boundary.
  \item \textbf{OWASP LLM Top 10 v1.1}
        \cite{owasp_llm_top10_v11}: LLM01 (prompt injection)
        and LLM06 (sensitive information disclosure) are the
        two entries our defenses address; the others are out
        of scope.
  \item \textbf{MCP specification}
        \cite{mcp_spec,anthropic_mcp_announce}: the spec
        defines the wire protocol; our spec-change
        recommendations are the minimum changes required
        to make the protocol-level isolation
        enforceable-by-default.
\end{itemize}

The defense blueprint is therefore not ad-hoc; it is the
intersection of the three standards-level priors applied
to the MCP protocol surface.